\chapter{Literature Review: ICT4D}
\begin{chapterintro}
This chapter gives discusses the related works in unemployment and Development.
\end{chapterintro}

\section{Development and Co-design}

Development is “the process of improving the quality of all human lives and capabilities by raising people’s levels of living, self-esteem, and freedom” \cite{todaro2015}. However, to improve the living standard/ “people’s levels of living” \cite{todaro2015} there is a need for money and other resources. To earn money or acquire resources, one needs to be employed or own a business. Various forms of technology are being increasingly seen as enablers of development \cite{npc2012}. However, development often advances with great imbalance, often termed -the digital divide \cite{feldmann2003, sambasivan2010, nicola2011}. In South Africa, the digital divide has been very great, causing a major disparity between the rich and the poor. As with the digital divide, there is also a great disparity in the unemployment rate in South Africa across different groupings. 

The South African unemployment level stands at 25.2\% in 2014 by the strict definition of unemployment. However, there is a great disparity in the unemployment levels when analysed along the groupings of race, age groups and skill levels. Zeroing into the unemployment level among the youth in South Africa (aged 15--24), the unemployment rate is higher, at 36.1\% and further zeroing into the black African youth in South Africa, the rate is even higher, at 39.4\% by 2014 \cite{statssa2015}. South Africa has the goal of slashing down unemployment to merely 6\% by 2030 \cite{npc2012}. These statistics indicate the nation’s need to pay greater attention to the unemployment issues among the youth and in the black communities in order to reach the 6\% mark. 

In South Africa, the financial services, information technology and business services have grown employment, however, mainly for the skilled \cite{npc2012}, thus a need for great efforts in the inclusion of the marginalised individuals. Technologies are generally considered unimportant in underdeveloped communities unless the technologies will contribute to the economic emancipation of the underdeveloped \cite{lorini2014}. 

South Africa is faced by a shortage of skilled labour, which raises the income of the skilled and thus increasing the wage inequality in the country \cite{npc2012}. The newly emergent field of Information Communication Technologies for Development (ICTD/ICT4D) has had several researchers investigation how to alleviate the inequalities: The Farming improvement project to provide agricultural information to farmers \cite{gandhi2007}, money transfer \cite{dyson2007}, job finding information \cite{qiang2012, chepken2012tele}, business starter software for under-developed start-ups \cite{brittan2011}, water service delivery \cite{ssozi2016} and many other infrastructural improvement and information access systems. The recent money transfer system MPESA which makes use of mobile technology to send money has been used by 14 million Kenyans, with transactions to the worth of 7 billion. These systems all aim to diminish the imbalances between rural and urban co-existent communities through the use of information and communication technologies. ICT4D/ICTD is a young discipline that focuses on ICT initiatives that contribute to social or economic development \cite{chepken2012trends}, usually targeted at poor regions of Africa, Asia and Latin America \cite{chepken2012trends}. ICT4D projects have been critiqued for being a cover-up for political agendas from those in power, thus ICTD projects should transparently state action plans, beneficiaries and benefactors \cite{elkhatib2015}. However, South African has identified the great potential in ICTs as listed in the national vision 2030 \cite{npc2012}. Although ICT4D has dated back to the 1990s, we understand it as an academic discipline as well as a field of practice in which Information and communication technologies are utilised to aim at development \cite{chepken2012trends}. In many past studies, ICT4D has been used to aid projects around health, education, agriculture and social communication \cite{david2013, chepken2012trends}. 

The Co-design approach to developing technology helps in improving people’s wellbeing, reducing creation costs (by engaging their very own ideas), and improving the adoption of technology \cite{steen2011}. Co-design can be perceived as a new term for Participatory Design or an improved version of participatory Design, its key principle is the involvement of all stakeholders in the design process to ensure that the end-products meet the users’ needs, are highly adapted and are sustainable \cite{david2013}. Co-design has become integrated with developmental initiatives because of two intersecting ideas: firstly in design, a move toward user involvement for product improvements and secondly, in development, a move towards social embeddedness of projects and the importance of locals in technology appropriation \cite{david2013}. The discussions around Co-design has five main aspects of consideration: stakeholders, context, ownership, social learning, and sustainability \cite{david2013}. In other words, Co-design is about selecting a representative set of stakeholders, performing design within end-users’ contexts, knowing artefact is owned/adopted, obtaining mutual learning and having ICTs that run on a long-term basis. 

\section{Causes of Unemployment}

Gitau attributes the unemployment problem to a lack of skills among the low-income \cite{gitau2012}. In addition, minimum wage and work conditions restrictions from bargaining councils and wage boards, which are applied uniformly irrespective of company size incur relatively higher labour costs on small businesses, \cite{kingdon2001}, which means that small business would not be motivated to employ many workers. This is true for the small subset of individuals who manage to enter the labour force, however, it is suspected that barriers such as disempowerment and restrictions of past apartheid period, lack of infrastructure, land and capital would bar unemployed persons from employment. Poverty discourages job-search activities, \cite{orkin1998} because searching for jobs incurs transport costs, CV creation costs and sometimes data costs. 

Meissner and Blake attribute unemployment to weaknesses in career guidance support that young people receive from the school and family \cite{meissner2011}. In their work, Meissner and Blake made use of personas to show that low-income community youth face challenges of family pressure to find jobs as soon as possible, low or no guidance on how to propel their careers as well as a need to manage the highly-constrained resource, time. The young individuals had to juggle between school and watching over younger ones, or between working and watching over younger ones. They proposed to address this problem by offering career guidance which involved using a computer and website, but also admitted that these individuals were time-constrained and lacked advisors who would give them advice when using the websites.  

Furthermore, research has shown that users with low education levels prefer to receive job notifications through their phones, via SMS or call as opposed to via emails \cite{white2012}. This indicates that they prefer simple mobile technology over computer-based or other advanced technologies. Individuals without college-level education tend to perform worse on computer-based solutions than those of college education level \cite{algarin2009}. However, email addresses are often requested or required on job searching apps and websites like Gumtree (gumtree.co.za), Indeed (indeed.co.za), and Freelancer (freelancer.co.za), thus, they are not directly suited to low-skilled persons. Thus, the proposed intervention in this dissertation should allow unemployed job seekers who are often low-skilled to use their accessible or preferred technology means. 

Moreover, Donner and Gitau identify an important category of technology users: mobile-only users \cite{donner2009}. Mobile-only users do not have the luxury of doing tasks with the most optimised devices (e.g. PCs for typing and video calls and mobile phones for making voice calls) because they have mobiles as their only means of communicating and task completion. These individuals could be better developed for, in a fashion that enables them to understand technology without the need to understand desktop computers. 

Meissner uncovered that users would like the best-paying jobs rather than the most fulfilling job, the findings also included that users do not get information about tertiary education and jobs frequently enough, and job availability information was sometimes outdated \cite{meissner2011}, highlighting the demand for relevant information and to access the information quickly \cite{bhavnani2008}. 

Since job seekers face so many challenges including the discussed: lack of skills, transport cost, CV creation cost, data cost, minimum wage and work conditions restrictions from bargaining councils and wage boards, and restrictions of past apartheid period, lack of infrastructure, land and capital, poverty and poor career guidance it is evident that that unemployment is a multidimensional issue which requires various agencies to tackle it. 

\section{ICT Interventions for Unemployment}

Several projects have attempted to use ICTs to tackle unemployment and poverty problems. Ummeli \cite{gitau2012} is a text-based mobile unemployment solution that was targeted at semi-literate individuals. We discuss this application because it was developed specifically for use in South Africa. Ummeli provides functionalities for its users to browse for jobs, post jobs and recommend jobs, create CVs and communicate with other job seekers. It was tested in the South African township of Khayelitsha \cite{gitau2012}. However, the application was registered for Vodacom cellular network users only. This meant that other unemployed candidates registered on other cellular networks such as MTN and Cell C would be left out. Babajob (babajob.com) is a job searching network that connects job seekers to employers, particularly connecting poor job seekers with few skills to elite employers. Babajob utilises cell phone numbers to register job seekers as an alternative to email addresses; and involves people, who may guide the use of technology by non-experts. Babajob’s interface is simple with very few steps required from a job applicant as opposed to Indeed and CareerJet job search platforms. Babajob also allows users to apply directly through the site, as well as subscribe to get SMS \cite{babajob2016}, however, the facilities are local to India.  

Like Babajob, Giraffe (giraffe.co.za), which is a job searching platform based in South Africa allows mobile numbers to be used to create an account. However, Giraffe has the shortcoming of requesting for South African identity numbers which then bars non-South Africans from accessing the job opportunities available through the website. 

Greataupair (greataupair.com/search/hire.cfm/careType/nanny) is another good website for job finding, which filters job seekers by age and salary, the website posts jobs and job seekers, however, it is designed specifically for housekeeping jobs. Greataupair requires an email to register an account. Like Babajob, the Kazi560 application in Kenya is an example of an application that successfully improves the lives of job seekers, job seekers can subscribe to SMS messages for a small fee per SMS \cite{qiang2012}.  

The Indian government announced the National Skills Development Policy, which aims at creating 500 million skilled employees by 2022, amongst its other goals. Thus, in Karnataka, India, the labour department implemented a national network of employment centres across the country where ‘candidates’/potential job seekers undergo an assessment of their background and skill-level and based on the outcome of the assessment, the centre identifies appropriate job vacancies or suitable training for the candidate. The centres are also reachable by calling in to speak to an agent or utilizing a government managed internet portal associated with the centres. Although the approach is sound, the state is unable to maintain the level of assistance required to reach their employability to the goal of creating 500 million skilled employees by 2022 and thus there is huge potential for technology solutions to improve the accessibility of employment services of the centres \cite{white2012}. Similarly the South African government has implemented the National Youth Development Agency (NYDA) to assist South Africans under 35 years with skills training, business start-up guidance and venture capital \cite{nyda2015}, yet like with Indian employment centres, the candidates reached through NYDA are limited, and huge numbers remain unemployed in the communities \cite{statssa2015}. 

Yet with the presence of several mobile interventions around unemployment, mobile Applications for Agricultural and Rural Development (m-ARDs) face challenges of finding funds to run. Several models of m-ARDs are therefore based on obtaining funding through SMS revenues which are low and often must depend on the external support of policymakers and government aid. Particularly in developing regions, m-ARDs face the challenges such as lack of local content, poor quality of available interfaces, influential powers due to monetary support, and because asymmetric access to information is often scarce in these regions \cite{qiang2012}. For example, employers having persistent access to internet and web platforms to post vacancies whereas job seekers are not always well empowered to post their profiles in order to market their skills. Voice-based systems can be utilised when working with the textually illiterate as has been successfully used in a rural setting in India to help individuals find employment \cite{white2012}. 

Some have proposed microtasks as an employment \cite{mtsweni2014, zyskowski2015}. Microtasks are small tasks often taken out of larger tasks and offer compensation to individuals from around the world \cite{mtsweni2014}. Examples of microwork include image location, item classification, translation and survey filling. Microtasks are offered on different microtask platforms accessible through smartphones, and other internet-connected devices. The following are examples of microtask platforms: Amazon’s Mechanical Turk (Mturk), InnoCentive \cite{innocentive2017}, Ajira, Jana and Samasource. However, research in a rural township in South Africa has shown that the community members mainly did not have computers and do not consider smartphones as a potential resource for income generation, as such ruling out their access to microtasks through smartphones \cite{lorini2014}. Microtask platforms offer job seekers non-traditional avenues to seek employment often with payment known upfront. Despite the great potential of microtasks to assist in unemployment, its accessibility is highly hampered because it is often presented in a digital form to be accessed with computers and internet connectivity which many low-skilled job seekers do not have \cite{gupta2012}. Mturk \cite{mturk2018} and TaskRabbit \cite{taskrabbit2013} are based in developed regions and restrict international workers from their use. Samasource obtains its tasks from the developed regions and makes them available in developing countries such as Kenya, South Africa and Uganda. Jana which started off as a small-scale microtask platform in Kenya and Rwanda \cite{eagle2009} has since been scaled-up for use in large organisations around the world, such as CNN, Unilever and United Nations. Jana is well suited to the developing regions in that it makes use of SMS which is cheap and USSD which does not incur any costs.

\section{Is Mobile for Development?}

In recent years, mobiles have become the most easily accessible and ubiquitous device for communication in constrained settings \cite{bhavnani2008, gsma2014}, and as such technologies should be built around the use of mobile devices for access to information and economic markets. Society has experienced a high recent rise in the use of mobile phones, in 2002 mobile subscriptions overtook fixed-lines \cite{feldmann2003}, mobile devices have many advantages which make them desirable. Mobiles are wireless, portable, have the affordance for intimate and personal control and are capable of supporting discrete digital tasks which can be interlaced \cite{donner2015}. Mobile devices are also useful for mobile money transfer \cite{donner2015}, however, the low-end devices/ non-programmable mobile phones do not allow for advanced USSD money transfer systems such as M-PESA \cite{chepken2012tele}. 

A study that involved job-seekers in Cape Town, Johannesburg, Nairobi and Windhoek indicated that although the job-seekers were able to use messaging and call services, they often did not have enough airtime on their devices, with an average of ZAR one in Cape town \cite{chepken2012tele}. Smart mobile devices are also useful for accessing the internet \cite{donner2013}, however, their usage does not override the need for computer-based internet access or Public Internet Access venues (PAVs) \cite{donner2013}. Wireless solutions have some advantages over their wired counterparts, they are faster, speedy to set up and also cost less; however, this is an underestimate when the cost for towers and solar power systems are included \cite{feldmann2003}. 

Intermittent (on and off) mobile data connection can be caused by device battery life, airtime balance, movement between places, changes in signal and congestion \cite{donner2015}. Finding a solution around all the possible causes of intermittent access would simply be infeasible but being knowledgeable about them can motivate the design of solutions that are mostly offline, battery conserving, and cost-effective, moreover, utilising airtime-less peer-to-peer communications (such as Bluetooth), and generally avoiding data intensive interactions.

\section{Including the Masses}

ICTs have great potential in disseminating information timelier and cost-effectively \cite{elijah2006}, however, acknowledging and designing with the knowledge that, communications are often braided \cite{densmore2013} over different channels is essential. The possibilities of channels in braided communication include SMS, Emailing, voice calls, video calls and social media, it is this important to develop solutions that cross platforms or better yet, solutions that are platform independent \cite{molapo2015} because developing an ICT that is suited for only one channel of communication or only a specific device can result in exclusion of people who frequent other channels and devices to communicate. 

The value that a user gains by being provided with new technologies depends on how well the technology fits into the user’s lifestyle (communicative ecology), in some instances new technologies may improve the lifestyle of its users when there is a small learning curve. This reinforces the need for simplicity as identified as a major HCI enabler as is reflective of Schneiderman’s 5th, 6th and 8th principle \cite{shneiderman2004}.

\section{Human Empowerment \& IT Sustainability}

Human empowerment is crucial for the sustainability of an intervention. Vancouver found that users will contribute to a system if they perceive intrinsic motivation such as having fun or helping people they feel affiliated with \cite{zimmerman2011}. In the Ummeli research, Gitau identified the most pressing needs of the participants was to locate and apply for jobs online, however, their going online was not exclusively for job finding but also for entertainment or entertainment-fostering purposes such as finding music, crafting profiles, sending birthday invitations, and befriending people \cite{gitau2012}. Co-design induces a sense of ownership as well as sustainability in a community \cite{david2013}, which is a result of the high engagement in every possible phase of participatory research. A sense of ownership of technology can be an empowerment realisation and sustainability reinforcement. However, in some cases, a technology system may require maintenance, in part by an IT expert or intermediary, and thus each developer should do a good job so that any developer to be able to maintain or build onto, in the case that the system is deemed valuable enough to be continually maintained. Heimerl and others found that it is important to empower a community to solve its own problem with existing knowledge and infrastructure \cite{heimerl2015}. 

Information dissemination is a crucial aspect of empowerment \cite{chepken2012tele}, for example, an artefact that informs people of vacancies and teaches them lessons useful for the work they would like to be employable for can be useful. Surfing the Internet is one way in which people can access information almost anytime, however, the ability to surf the internet is inhibited by the lack of know-how \cite{lorini2014}, mobile access (small sized devices and gesture inconveniences) and search engine results and presentation format. In a recent study with individuals who were semi-skilled and preferably looking for skills training, 60\% of the individuals indicated that they did not find satisfactory information, as search engines did not offer relevant information, did not limit the results to information within their geographical location and gave an overwhelming number of results when they were looking for training courses \cite{paihama2016}. Obtaining ten results was overwhelming for the participants, however, when the number of results was reduced to four, the participants were able to complete their tasks. Because of the shortcomings of search engines and search engines users’ inabilities to search efficiently, word of mouth or person to person communication can be more relevant for finding job opportunities and skills training services. 

A system can be sustainable either by self-maintaining or being human maintained. Having a self-maintaining IT system is close to impossible, in a society where information changes in unpredicted ways whereas human maintenance is possible and feasible. Ensuring that users derive value from the maintenance of a system can motivate them to keep maintaining it. \cite{schwartz2013} found that the unintended uses of devices may motivate users to maintain devices, to learn advanced tasks and possibly instil a sense of ownership in users. A sense of ownership can reinforce freedom and lead to empowerment and sustainability. However, when the empowered individuals are not IT specialists a system may need to be maintained by a specialised IT expert, and thus the system would not be completely sustained by the users, however, if the users are empowered enough they could deem the system valuable enough to invest in getting an external party to maintain it. 

Sustainability can also be approached by ensuring the design of an artefact is not only with affordances for functionalities but is also afforded to let users know of its affordances \cite{mcgrenere2000} as opposed to, having hidden affordances. Designs should be useful for new users, easy to maintain and documented in a way that they can be built unto.
